% -*- coding: utf-8 -*-
%-------------------------designed by zcf--------------
\documentclass[UTF8,a4paper,10pt]{ctexart}
\usepackage[left=3.17cm, right=3.17cm, top=2.74cm, bottom=2.74cm]{geometry}
\usepackage{amsmath}
\usepackage{graphicx,subfig}
\usepackage{float}
\usepackage{cite}
\usepackage{caption}
\usepackage{enumerate}
\usepackage{booktabs} %表格
\usepackage{multirow}
\newcommand{\tabincell}[2]{\begin{tabular}{@{}#1@{}}#2\end{tabular}}  %表格强制换行
%-------------------------字体设置--------------
% \usepackage{times} 
% \usepackage{ctex} % ctexart 已经加载 ctex，避免重复加载
\setCJKmainfont[ItalicFont=Noto Sans CJK SC Bold, BoldFont=Noto Serif CJK SC Black]{Noto Serif CJK SC}
\newcommand{\yihao}{\fontsize{26pt}{36pt}\selectfont}           % 一号, 1.4 倍行距
\newcommand{\erhao}{\fontsize{22pt}{28pt}\selectfont}          % 二号, 1.25倍行距
\newcommand{\xiaoer}{\fontsize{18pt}{18pt}\selectfont}          % 小二, 单倍行距
\newcommand{\sanhao}{\fontsize{16pt}{24pt}\selectfont}  %三号字
\newcommand{\xiaosan}{\fontsize{15pt}{22pt}\selectfont}        % 小三, 1.5倍行距
\newcommand{\sihao}{\fontsize{14pt}{21pt}\selectfont}            % 四号, 1.5 倍行距
\newcommand{\banxiaosi}{\fontsize{13pt}{19.5pt}\selectfont}    % 半小四, 1.5倍行距
\newcommand{\xiaosi}{\fontsize{12pt}{18pt}\selectfont}            % 小四, 1.5倍行距
\newcommand{\dawuhao}{\fontsize{11pt}{11pt}\selectfont}       % 大五号, 单倍行距
\newcommand{\wuhao}{\fontsize{10.5pt}{15.75pt}\selectfont}    % 五号, 单倍行距
%-------------------------章节名----------------
% \usepackage{ctexcap} % ctexart 已经加载标题设置，避免与 ctexart 冲突 
\ctexset{
  section = {
    name = {,、},
    number = \chinese{section}
  },
  subsection = {
    name = {（,）},
    number = \chinese{subsection}
  },
  subsubsection = {
    name = {,.},
    number = \arabic{subsubsection}
  }
}
%-------------------------页眉页脚--------------
\usepackage{fancyhdr}
\pagestyle{fancy}
\lhead{\kaishu \leftmark}
% \chead{}
\rhead{\kaishu 并行程序设计实验报告}%加粗\bfseries 
\lfoot{}
\cfoot{\thepage}
\rfoot{}
\renewcommand{\headrulewidth}{0.1pt}  
\renewcommand{\footrulewidth}{0pt}%去掉横线
\newcommand{\HRule}{\rule{\linewidth}{0.5mm}}%标题横线
\newcommand{\HRulegrossa}{\rule{\linewidth}{1.2mm}}
%-----------------------伪代码------------------
\usepackage{algorithm}  
\usepackage{algorithmicx}  
\usepackage{algpseudocode}  
\floatname{algorithm}{Algorithm}  
\renewcommand{\algorithmicrequire}{\textbf{Input:}}  
\renewcommand{\algorithmicensure}{\textbf{Output:}} 
\usepackage{lipsum}  
\makeatletter
\newenvironment{breakablealgorithm}
  {% \begin{breakablealgorithm}
  \begin{center}
     \refstepcounter{algorithm}% New algorithm
     \hrule height.8pt depth0pt \kern2pt% \@fs@pre for \@fs@ruled
     \renewcommand{\caption}[2][\relax]{% Make a new \caption
      {\raggedright\textbf{\ALG@name~\thealgorithm} ##2\par}%
      \ifx\relax##1\relax % #1 is \relax
         \addcontentsline{loa}{algorithm}{\protect\numberline{\thealgorithm}##2}%
      \else % #1 is not \relax
         \addcontentsline{loa}{algorithm}{\protect\numberline{\thealgorithm}##1}%
      \fi
      \kern2pt\hrule\kern2pt
     }
  }{% \end{breakablealgorithm}
     \kern2pt\hrule\relax% \@fs@post for \@fs@ruled
  \end{center}
  }
\makeatother
%------------------------代码-------------------
\usepackage{xcolor} 
\usepackage{listings} 
\lstset{ 
breaklines,%自动换行
basicstyle=\small,
% escapeinside=``, % 已关闭：否则代码注释中 CRT_MOD_COUNT 这类下划线会触发 Missing $ inserted
keywordstyle=\color{ blue!70} \bfseries,
commentstyle=\color{red!50!green!50!blue!50},% 
stringstyle=\ttfamily,% 
extendedchars=false,% 
linewidth=\textwidth,% 
numbers=left,% 
numberstyle=\tiny \color{blue!50},% 
frame=trbl% 
rulesepcolor= \color{ red!20!green!20!blue!20} 
}
%------------超链接----------
\usepackage[colorlinks,linkcolor=black,anchorcolor=blue]{hyperref}
%------------------------TODO-------------------
\usepackage{enumitem,amssymb}
\newlist{todolist}{itemize}{2}
\setlist[todolist]{label=$\square$}
% for check symbol 
\usepackage{pifont}
\newcommand{\cmark}{\ding{51}}%
\newcommand{\xmark}{\ding{55}}%
\newcommand{\done}{\rlap{$\square$}{\raisebox{2pt}{\large\hspace{1pt}\cmark}}\hspace{-2.5pt}}
\newcommand{\wontfix}{\rlap{$\square$}{\large\hspace{1pt}\xmark}}
%------------------------水印-------------------
\usepackage{tikz}
\usepackage{xcolor}
\usepackage{eso-pic}

\newcommand{\watermark}[3]{\AddToShipoutPictureBG{
\parbox[b][\paperheight]{\paperwidth}{
\vfill%
\centering%
\tikz[remember picture, overlay]%
  \node [rotate = #1, scale = #2] at (current page.center)%
    {\textcolor{gray!80!cyan!30!magenta!30}{#3}};
\vfill}}}



%———————————————————————————————————————————正文———————————————————————————————————————————————
%----------------------------------------------
\begin{document}
\begin{titlepage}
    \begin{center}
    \includegraphics[width=0.8\textwidth]{NKU.png}\\[1cm]    
    \textsc{\Huge \kaishu{\textbf{南\ \ \ \ \ \ 开\ \ \ \ \ \ 大\ \ \ \ \ \ 学}} }\\[0.9cm]
    \textsc{\huge \kaishu{\textbf{计\ \ 算\ \ 机\ \ 学\ \ 院}}}\\[0.5cm]
    \textsc{\Large \textbf{并行程序设计期末实验报告}}\\[0.8cm]
    \HRule \\[0.9cm]
    { \LARGE \bfseries 基于 Pthread、OpenMP 与 CRT 的\\[0.3cm] NTT 多项式乘法并行优化}\\[0.4cm]
    \HRule \\[2.0cm]
    \centering
    \textsc{\LARGE \kaishu{姓名\ :\ \underline{\makebox[5em][c]{TODO}}}}\\[0.5cm]
    \textsc{\LARGE \kaishu{学号\ :\ \underline{\makebox[8em][c]{TODO}}}}\\[0.5cm]
    \textsc{\LARGE \kaishu{年级\ :\ \underline{\makebox[5em][c]{TODO}}}}\\[0.5cm]
    \textsc{\LARGE \kaishu{专业\ :\ \underline{\makebox[8em][c]{计算机科学与技术}}}}\\[0.5cm]
    \textsc{\LARGE \kaishu{指导教师\ :\ \underline{\makebox[5em][c]{TODO}}}}\\[0.5cm]
    \vfill
    {\Large \today}
    \end{center}
\end{titlepage}
%-------------摘------要--------------
\newpage
\thispagestyle{empty}
\renewcommand{\abstractname}{\kaishu \sihao \textbf{摘要}}
    \begin{abstract}

        \noindent  %顶格
        本实验在数论变换（Number Theoretic Transform, NTT）多项式乘法的迭代实现基础上，
        分别使用 Pthread 与 OpenMP 完成了多线程并行化，
        并进一步针对大模数 $p=1337006139375617$ 实现了基于中国剩余定理（CRT）的进阶优化。
        报告以 32 位标量迭代 NTT（V0）作为性能基准，按版本演进的方式讨论了 OpenMP（V1）
        与 Pthread（V2--V6）多种任务划分策略，重点对比了按 block 划分、按 butterfly 展平划分、
        连续 flat 划分、whole-pipeline 线程池以及 cyclic-flat 在负载均衡、同步开销
        与访存局部性上的差异。在三组 32 位模数实验中，OpenMP 8 线程平均延迟为
        $21.7870$\,ms，Pthread whole-pipeline 8 线程平均为 $23.2387$\,ms，OpenMP 较 Pthread
        快约 $6.7\%$。在大模数场景下，直接基于 \texttt{\_\_int128} 的标量 NTT 耗时
        $109.006$\,ms，OpenMP 版本降至 $45.704$\,ms，进一步采用 4 模数分组 CRT 结合
        Montgomery NTT 与 Fast Garner 合并的版本（V10-2）在 8 线程下三次复跑平均
        $33.6254$\,ms，相对直接大模数标量加速约 $3.24\times$。

        \vspace{0.5em}

        \noindent\textbf{关键字：}NTT；多项式乘法；Pthread；OpenMP；中国剩余定理；Montgomery 模乘；Garner 合并
\end{abstract}
%----------------------------------------------------------------
\tableofcontents
%----------------------------------------------------------------
\newpage
\watermark{60}{10}{NKU}
\setcounter{page}{1}

%================================================================
\section{实验概述}
%================================================================

\subsection{实验目的}

本次实验对应课程 Lab3 多线程编程作业，要求在已有串行算法的基础上，分别使用 POSIX Threads（Pthread）\cite{butenhof1997programming}
和 OpenMP\cite{dagum1998openmp} 完成所选选题的并行化实现，并系统分析不同任务划分策略、线程数与同步机制对性能的影响。
本人选择的选题为 NTT 多项式乘法优化，主要目标是：

\begin{enumerate}
  \item 基于实验框架完成迭代 NTT 的串行基线实现，作为后续并行化的参考；
  \item 分别使用 OpenMP 与 Pthread 实现 NTT 多线程版本，比较任务划分、同步策略
        和线程管理开销对性能的影响；
  \item 在三组 32 位 NTT-friendly 模数上测试 $1/2/4/8$ 线程的性能，分析扩展性、
        加速比与并行效率；
  \item 完成进阶部分：针对大模数 $p = 1337006139375617$ 实现基于中国剩余定理\cite{knuth1997art}
        的多模数并行优化，对比直接大模数 NTT 与 CRT 大模数 NTT 的性能；
  \item 形成对“线程创建/销毁开销”、“后期 stage 负载不均”、“访存局部性”
        等典型并行性能问题的工程理解。
\end{enumerate}

\subsection{选题背景}

多项式乘法是密码学、信号处理与同态加密等领域的基础运算。
在同态加密与基于 ideal lattice 的密码学方案中，模数通常较大，
朴素 $O(n^2)$ 实现完全不可接受。学术与工程实现普遍采用 NTT
将其转化为 $O(n \log n)$\cite{cooley1965algorithm}。
NTT 在数论层面替代了 FFT 的复数单位根，避免了浮点误差，并便于在
有限域 $\mathbb{Z}_p$ 上进行精确计算，被广泛应用于 ideal lattice
密码原语的高效实现\cite{longa2016speeding,harvey2014faster}。

NTT 的迭代实现由若干 stage 组成，每层 stage 内的 butterfly 运算之间相互独立，
是天然的并行计算单元；但不同 stage 之间存在数据依赖，必须借助同步原语
确保上一层全部完成后才能进入下一层。这种“同步等价于全局 barrier”的模式
使其成为研究多线程任务划分、负载均衡和同步开销的良好对象，与本课程
Lab3 中关于 Pthread 静态线程与 OpenMP 隐式 barrier 的讨论高度契合。

\subsection{本文完成内容}

依据 Lab3 Pthread 与 OpenMP 实验指导书以及 NTT 选题介绍的要求，本文完成的主要工作包括：

\begin{itemize}
  \item 实现 32 位 NTT 标量基线 V0，采用 bit-reversal 与 Montgomery 模乘\cite{montgomery1985modular}，
        作为后续并行化的对照；
  \item 使用 OpenMP 实现 stage 内 block 并行的多线程 NTT（V1），通过外层 \texttt{parallel}
        区域、\texttt{single} 预计算单位根和 \texttt{for schedule(static)} 完成并行；
  \item 使用 Pthread 完成五个版本的并行化：block-parallel（V2）、flat-butterfly（V3）、
        连续 flat（V4）、whole-pipeline 线程池（V5）、whole-pipeline+cyclic-flat（V6），
        并对各版本性能进行对比分析；
  \item 完成 $1/2/4/8$ 线程数下 OpenMP 与 Pthread 的扩展性实验；
  \item 实现大模数 NTT 的标量版（V7）与 OpenMP 版（V8），使用 \texttt{\_\_int128}
        处理 60 位以内模数；
  \item 完成 CRT 大模数进阶优化：5 模数初版（V9）、4 模数 + Montgomery（V10-1）、
        4 模数分组 CRT + Montgomery NTT + Fast Garner（V10-2）；
  \item 撰写实验报告，对所有版本的性能、加速比、并行效率与实现复杂度进行系统比较。
\end{itemize}

需要说明的是，按照 Lab3 实验要求，本文重点为 Pthread 与 OpenMP 多线程优化，
NTT 选题介绍中提到的“四分 NTT”优化未在本实验中实现，SIMD/NEON 仅作为上一实验的
背景出现，在本实验中以非 SIMD 的标量 NTT 为基线出发，集中讨论多线程并行化与 CRT
大模数优化。

%================================================================
\section{算法原理}
%================================================================

\subsection{多项式乘法定义}

给定两个多项式
\[
  f(x) = a_{n-1}x^{n-1} + \cdots + a_1 x + a_0,\qquad
  g(x) = b_{m-1}x^{m-1} + \cdots + b_1 x + b_0,
\]
其乘积 $h(x) = f(x) g(x)$ 的 $x^k$ 项系数为
\[
  h_k = \sum_{i=0}^{k} a_i \, b_{k-i}.
\]
朴素卷积实现的时间复杂度为 $O(nm)$，对于 $n, m$ 量级达到 $10^5$ 的输入不具备实用性。

\subsection{NTT 与有限域上的单位根}

设 $p$ 为素数且 $L \mid (p-1)$，$g$ 为 $p$ 的一个原根，定义 $\mathbb{Z}_p$ 上的 $L$ 次单位根
\[
  \omega_L = g^{(p-1)/L} \bmod p.
\]
$\omega_L$ 具备与复数单位根相同的代数性质，因此可以将 FFT 中所有复数单位根换为 $\omega_L$，
将所有复数运算换为模 $p$ 运算，即得到 NTT。多项式乘法的标准 NTT 流程为：

\begin{enumerate}
  \item 将 $f$、$g$ 的系数向量补零至 $L = 2^{\lceil \log_2 (n+m-1) \rceil}$；
  \item 分别对两向量做正向 NTT，得到点值表示；
  \item 点值逐点相乘，得到 $h$ 的点值表示；
  \item 对结果做逆 NTT，并乘以 $L^{-1}$，得到 $h$ 的系数表示；
  \item 截取前 $n + m - 1$ 项作为最终结果。
\end{enumerate}

\subsection{迭代 NTT 与 butterfly 分层}

迭代 NTT 首先对输入做 bit-reversal 置换，随后按层（stage）执行 butterfly。
单次正向 NTT 的伪代码如下：

\begin{breakablealgorithm} 
  \caption{迭代 NTT 主体——按 stage 执行 butterfly}
  \begin{algorithmic}[1]
      \Require 长度 $L = 2^k$ 的系数向量 $a$，模数 $p$，原根 $g$
      \Ensure $a$ 的 NTT 变换结果（原地）
      \State 对 $a$ 应用 bit-reversal 置换
      \For{$b = 1, 2, 4, \dots, L/2$}
          \State $B \gets 2b$
          \State $\omega_n \gets g^{(p-1)/B} \bmod p$
          \State 预计算 $\mathrm{stageRoots}[j] = \omega_n^j \bmod p$，$j = 0, 1, \dots, b-1$
          \For{$k = 0, B, 2B, \dots, L - B$}
              \For{$j = 0, 1, \dots, b - 1$}
                  \State $u \gets a[k+j]$
                  \State $t \gets a[k+j+b] \cdot \mathrm{stageRoots}[j] \bmod p$
                  \State $a[k+j] \gets (u + t) \bmod p$
                  \State $a[k+j+b] \gets (u - t) \bmod p$
              \EndFor
          \EndFor
      \EndFor
      \State \Return $a$
  \end{algorithmic}  
\end{breakablealgorithm}

同一 stage 内每个 block 与每个 butterfly 都是独立运算，天然支持多线程；
不同 stage 之间存在写后读依赖，必须通过 barrier 或 OpenMP 的隐式同步保证一致性。

\subsection{Montgomery 模乘}

对于 32 位 NTT-friendly 模数，模乘出现在每个 butterfly 中，是 NTT 的主要计算量来源。
Montgomery 模乘\cite{montgomery1985modular} 选取参数 $R = 2^{32}$，
通过预计算 $\mu = -p^{-1} \bmod R$，将原本的 \texttt{\%} 运算替换为乘法 + 加法 + 移位组合：
\[
  \mathrm{redc}(T) = \bigl(T + ((T \bmod R) \cdot \mu \bmod R) \cdot p\bigr) / R.
\]
这样 \texttt{ab \% p} 的代价显著低于直接取模，并且在不依赖 SIMD 的标量代码中也能取得收益。

\subsection{中国剩余定理与 Garner 合并}

当目标模数 $P$ 超出 32 位甚至接近 60 位时，单次模乘必须使用 \texttt{\_\_int128} 或
更高精度类型，开销显著增加。CRT 多模数策略的核心思想是：

\begin{enumerate}
  \item 选取若干互素的 32 位 NTT-friendly 模数 $p_0, p_1, \dots, p_{k-1}$，
        其乘积 $\prod p_i$ 必须严格大于卷积单点最大值；
  \item 分别在每个小模数下完成系数 $\bmod\, p_i$ 的 NTT 多项式乘法；
  \item 利用 Garner 算法\cite{knuth1997art} 将各小模数下的结果重构出 $\bmod\, P$ 的结果。
\end{enumerate}

Garner 合并的关键是预计算各对 $p_i, p_j$ 的模逆元 $\mathrm{inv}_{ij} = p_i^{-1} \bmod p_j$，
使得合并阶段只剩下若干次乘加而非昂贵的大数除法，显著降低 CRT 合并阶段的常数。

%================================================================
\section{串行基线与数据结构}
%================================================================

\subsection{V0：32 位标量迭代 NTT 基线}

V0 是本文所有并行版本的对照基准。它是非 SIMD 的纯标量实现，采用 Montgomery 表示存储
stageRoots，避免每个 butterfly 重新进入 \texttt{\%} 运算。其关键数据结构包括：

\begin{itemize}
  \item \texttt{firstTransformed}, \texttt{secondTransformed}：两个长度 $L \le 2^{18}$
        的静态数组，作为输入多项式补零后存放系数的 buffer；
  \item \texttt{stageRoots[]}, \texttt{stageRootsMontgomery[]}：每个 stage 预计算的
        单位根表，Montgomery 形式存储以便快速模乘；
  \item \texttt{bitReverseTable[]}：bit-reversal 置换表，在每次 NTT 开始前应用。
\end{itemize}

V0 在三组 32 位输入下的延迟约 $44\sim46$\,ms，作为后续 OpenMP 与 Pthread 1 线程对照基线。

\subsection{正确性验证}

实验框架提供两个正确性检查函数：

\begin{itemize}
  \item \texttt{fCheck}：用于前三个 32 位模数的输入；
  \item \texttt{fCheckLarge}：用于第四个大模数 $p = 1337006139375617$ 的输入。
\end{itemize}

本实验中所有列入正式结果的版本均通过对应检查并输出``多项式乘法结果正确''或
``大模数多项式乘法结果正确''。测试脚本在结果输出之后向日志路径写入时偶尔
出现权限错误，但发生在算法完成之后，不影响算法正确性，本文不展开。

%================================================================
\section{OpenMP 并行设计}
%================================================================

\subsection{设计思路}

按照实验指导书的提示，OpenMP 实现采用以下设计：

\begin{enumerate}
  \item 将 \texttt{\#pragma omp parallel} 上提到 NTT 主循环之外，让线程组在
        stage 循环中复用，避免每个 stage 都进行 fork-join 带来的额外开销；
  \item 由 \texttt{\#pragma omp single} 执行每层 stageRoots 的预计算；
  \item 使用 \texttt{\#pragma omp for schedule(static)} 并行处理当前 stage
        的所有 block；
  \item 利用 \texttt{single} 与 \texttt{for} 的隐式 barrier 保证``单位根计算完毕''
        与``当前 stage 全部完成''两个时间点；
  \item 同时对清零、输入拷贝、点值乘法、逆 NTT 后的缩放以及最终结果回拷
        都使用 \texttt{\#pragma omp for} 进行并行化。
\end{enumerate}

\subsection{关键代码片段}

\begin{lstlisting}[title=OpenMP NTT 主体伪代码,frame=trbl,language={C++}]
#pragma omp parallel num_threads(threadCount)
{
    // 1) bit-reversal `置换并行`
    #pragma omp for schedule(static)
    for (int i = 0; i < L; ++i) {
        int j = bitReverseTable[i];
        if (j > i) std::swap(a[i], a[j]);
    }

    // 2) stage `循环`
    for (int b = 1; b < L; b <<= 1) {
        int B = b << 1;

        // `由一个线程预计算本层单位根；隐式 barrier 保证可见性`
        #pragma omp single
        {
            uint32_t w  = mont_one;
            uint32_t wn = mont_pow(omega, (P - 1) / B);
            for (int j = 0; j < b; ++j) {
                stageRootsMontgomery[j] = w;
                w = mont_mul(w, wn);
            }
        }

        // block `级并行：每个线程负责若干个 block`
        #pragma omp for schedule(static)
        for (int k = 0; k < L; k += B) {
            for (int j = 0; j < b; ++j) {
                uint32_t u = a[k + j];
                uint32_t t = mont_mul(a[k + j + b],
                                      stageRootsMontgomery[j]);
                a[k + j]     = add_mod(u, t);
                a[k + j + b] = sub_mod(u, t);
            }
        }
        // `隐式 barrier 保证下一层之前本层完成`
    }
}
\end{lstlisting}

\subsection{OpenMP 与 V0 三组 32 位输入对比}

\begin{table}[H]
  \centering
  \begin{tabular}{lccc}
  \toprule  
  \textbf{输入} & \textbf{V0 标量} (ms) & \textbf{V1 OpenMP 8 线程} (ms) & \textbf{加速比}\\
  \midrule
  input 1, $p = 7340033$    & 44.4277 & 22.5479 & 约 $1.97\times$ \\
  input 2, $p = 104857601$  & 45.3791 & 23.0579 & 约 $1.97\times$ \\
  input 3, $p = 469762049$  & 45.8223 & 23.5296 & 约 $1.95\times$ \\
  \bottomrule
  \end{tabular}
  \caption{V0 标量与 V1 OpenMP 8 线程在三组 32 位输入上的对比}
  \label{tab:v0vsomp}
\end{table}

可以看到，OpenMP 8 线程在三组输入上均能稳定取得近 $2\times$ 加速。注意这里以单线程的
V0 标量版作为对照，已经包含了 Montgomery 模乘等标量优化，因此加速比受到串行段
（bit-reversal、单位根预计算）的较大牵制。下一节将以 OpenMP 自身 1 线程为基准重新
计算并行效率。

\subsection{OpenMP 线程数扩展性实验}

\begin{table}[H]
  \centering
  \begin{tabular}{ccccc}
  \toprule  
  \textbf{线程数} & $p=7340033$ & $p=104857601$ & $p=469762049$ & \textbf{三组平均}\\
  \midrule
  1 & 45.0691 & 45.9454 & 46.2821 & 45.7655 \\
  2 & 30.4181 & 31.1385 & 31.3640 & 30.9735 \\
  4 & 24.3344 & 24.4606 & 24.7056 & 24.5002 \\
  8 & 21.3026 & 21.9200 & 22.1383 & \textbf{21.7870} \\
  \bottomrule
  \end{tabular}
  \caption{OpenMP 在 1/2/4/8 线程下三组输入的运行时间（单位：ms）}
  \label{tab:omp-scale}
\end{table}

\begin{table}[H]
  \centering
  \begin{tabular}{ccc}
  \toprule  
  \textbf{线程数} & \textbf{平均加速比} & \textbf{平均并行效率}\\
  \midrule
  1 & $1.000\times$ & 100.0\% \\
  2 & $1.478\times$ & 73.9\%  \\
  4 & $1.868\times$ & 46.7\%  \\
  8 & $2.101\times$ & 26.3\%  \\
  \bottomrule
  \end{tabular}
  \caption{OpenMP 平均加速比与并行效率（以 OpenMP 1 线程为基准）}
  \label{tab:omp-speedup}
\end{table}

OpenMP 随着线程数增加持续变快，8 线程时取得最佳平均延迟 $21.79$\,ms。
但加速比明显低于线性增长，主要原因有：

\begin{itemize}
  \item bit-reversal 阶段虽然可以并行，但写冲突要求按 $j > i$ 判断，引入分支开销；
  \item 每层单位根预计算由一个线程完成（\texttt{single}），是串行段；
  \item 每层 \texttt{for} 结束的隐式 barrier 同步开销随线程数增加；
  \item 后期 stage $b$ 变大、block 数变小，并行粒度下降；
  \item NTT 是访存密集型计算，线程数增加后内存带宽成为瓶颈。
\end{itemize}

%================================================================
\section{Pthread 并行设计与优化演进}
%================================================================

\subsection{V2：block-parallel}

V2 是 Pthread 入门版本，思路与 OpenMP V1 类似：对每个 stage，按 block 把任务分给所有
线程。但每次 NTT 调用结束都需要 \texttt{pthread\_create} 与 \texttt{pthread\_join}，
开销显著。三组 32 位输入的延迟为：

\begin{itemize}
  \item $p = 7340033$: $38.3756$\,ms
  \item $p = 104857601$: $41.2726$\,ms
  \item $p = 469762049$: $40.4079$\,ms
\end{itemize}

主要瓶颈在于后期 stage 的 \texttt{blockCount} 急剧减少：当 $b = L/2$ 时只有 $1$ 个 block，
实际上只有 $1$ 个线程在工作，其余 $7$ 个线程闲置等待 barrier。

\subsection{V3：flat-butterfly}

V3 将每层的 $L/2$ 个 butterfly 拉平为线性任务序列，按线程 id 做 cyclic 或区间划分，
使每一层都能让所有线程均匀承担工作。三组 32 位输入：

\begin{itemize}
  \item $p = 7340033$: $35.1163$\,ms
  \item $p = 104857601$: $35.6363$\,ms
  \item $p = 469762049$: $36.0163$\,ms
\end{itemize}

性能从 V2 的 $38$--$41$\,ms 下降到 $35$--$36$\,ms，直接验证了后期 stage 闲置是
V2 的主要瓶颈，这一改造正是 Lab3 实验指导书所强调的``合理任务划分''。

\subsection{V4：连续 flat-butterfly}

V4 在 V3 的基础上，将 cyclic 划分换为连续区间划分，试图改善 cache 局部性。结果如下：

\begin{itemize}
  \item $p = 7340033$: $36.0632$\,ms
  \item $p = 104857601$: $38.7218$\,ms
  \item $p = 469762049$: $35.3413$\,ms
\end{itemize}

V4 在不同输入上波动较大，与 V3 相比未取得稳定优势。原因是 NTT 中 butterfly 的访存模式
由 stage 控制，其访存对的间隔依赖 $b$，连续划分并不一定能映射到连续的内存块上。

\subsection{V5：whole-pipeline 线程池}

V5 是本文最终采用的 Pthread 版本。它将整个多项式乘法流水线（清零、输入拷贝、两次正向
NTT、点值乘法、逆 NTT、缩放、结果拷贝）都纳入同一组 Pthread 线程的工作中，
整个流水线只创建一次线程组。

线程 0 负责 bit-reversal 与每层 stageRoots 预计算，其余线程在 barrier 处等待；
当线程 0 完成预计算后释放 barrier，所有线程并行执行该 stage 的 butterfly；
本 stage 完成后再以 barrier 同步进入下一阶段。三组 32 位输入的延迟为：

\begin{itemize}
  \item $p = 7340033$: $25.9370$\,ms
  \item $p = 104857601$: $26.5954$\,ms
  \item $p = 469762049$: $26.8529$\,ms
\end{itemize}

相对 V2--V4 大幅下降，证明``线程创建与销毁开销''是基础 Pthread 实现的主要瓶颈之一，
这也正是 Lab3 指导书中``Pthread 编程范式''一节反复强调的静态线程优势\cite{butenhof1997programming}。

\subsection{V6：cyclic-flat 反例}

V6 在 V5 的 whole-pipeline 框架上把 block 划分换成 cyclic flat-butterfly 划分。结果反而退化：

\begin{itemize}
  \item $p = 7340033$: $40.6159$\,ms
  \item $p = 104857601$: $41.2744$\,ms
  \item $p = 469762049$: $42.4182$\,ms
\end{itemize}

原因主要有：

\begin{enumerate}
  \item flat-butterfly 每个 butterfly 都需要从全局任务 id 反推 \texttt{blockIndex}
        和 \texttt{offset}，引入整数除法与取模；
  \item cyclic 分配破坏了同一线程访问连续 block 的局部性，每次 butterfly 都跳跃到
        不同的 cache line；
  \item 在 V5 的 whole-pipeline 框架下，原本 block 级并行已经能让 8 个线程在 NTT 的
        大部分 stage 中保持忙碌，进一步细化任务划分带来的负载均衡收益小于额外开销。
\end{enumerate}

\subsection{Pthread 版本演进对比}

\begin{table}[H]
  \centering
  \begin{tabular}{lccc}
  \toprule  
  \textbf{版本} & $p=7340033$ & $p=104857601$ & $p=469762049$ \\
  \midrule
  V2 block-parallel             & 38.3756          & 41.2726          & 40.4079          \\
  V3 flat-butterfly             & 35.1163          & 35.6363          & 36.0163          \\
  V4 contiguous flat-butterfly  & 36.0632          & 38.7218          & 35.3413          \\
  V5 whole-pipeline pool        & \textbf{25.9370} & \textbf{26.5954} & \textbf{26.8529} \\
  V6 whole-pipeline cyclic-flat & 40.6159          & 41.2744          & 42.4182          \\
  \bottomrule
  \end{tabular}
  \caption{Pthread 各版本三组 32 位输入运行时间对比（单位：ms）}
  \label{tab:pthread-versions}
\end{table}

\subsection{Pthread V5 关键代码片段}

\begin{lstlisting}[title=Pthread whole-pipeline 线程函数主体（简化）,frame=trbl,language={C++}]
void* worker(void* arg) {
    int tid = ((WorkerArg*)arg)->tid;

    // `阶段 1：参与清零 / 输入拷贝`
    do_clear_and_copy(tid);
    pthread_barrier_wait(&stage_barrier);

    // `阶段 2：正向 NTT（first）`
    forward_ntt_pipeline(firstTransformed, tid);

    // `阶段 3：正向 NTT（second）`
    forward_ntt_pipeline(secondTransformed, tid);

    // `阶段 4：点值乘法`
    pointwise_multiply(tid);
    pthread_barrier_wait(&stage_barrier);

    // `阶段 5：逆 NTT`
    inverse_ntt_pipeline(firstTransformed, tid);

    // `阶段 6：缩放 + 结果回拷`
    scale_and_writeback(tid);
    pthread_barrier_wait(&stage_barrier);
    return nullptr;
}

void forward_ntt_pipeline(uint32_t* a, int tid) {
    // bit-reversal `由线程 0 完成`
    if (tid == 0) apply_bit_reversal(a);
    pthread_barrier_wait(&stage_barrier);

    for (int b = 1; b < L; b <<= 1) {
        if (tid == 0) precompute_stage_roots(b);
        pthread_barrier_wait(&stage_barrier);

        int B = b << 1;
        // `按 block 做静态划分`
        for (int k = tid * B; k < L; k += B * THREAD_COUNT) {
            butterfly_block(a, k, b);
        }
        pthread_barrier_wait(&stage_barrier);
    }
}
\end{lstlisting}

\subsection{Pthread 线程数扩展性实验（基于 V5）}

\begin{table}[H]
  \centering
  \begin{tabular}{ccccc}
  \toprule  
  \textbf{线程数} & $p=7340033$ & $p=104857601$ & $p=469762049$ & \textbf{三组平均}\\
  \midrule
  1 & 45.0159 & 45.7279 & 46.0507 & 45.5982 \\
  2 & 31.1355 & 31.7720 & 33.8902 & 32.2659 \\
  4 & 24.0226 & 25.4599 & 24.6986 & 24.7270 \\
  8 & 22.5886 & 23.6726 & 23.4548 & \textbf{23.2387} \\
  \bottomrule
  \end{tabular}
  \caption{Pthread V5 在 1/2/4/8 线程下三组输入的运行时间（单位：ms）}
  \label{tab:pthread-scale}
\end{table}

\begin{table}[H]
  \centering
  \begin{tabular}{ccc}
  \toprule  
  \textbf{线程数} & \textbf{平均加速比} & \textbf{平均并行效率}\\
  \midrule
  1 & $1.000\times$ & 100.0\% \\
  2 & $1.415\times$ & 70.7\%  \\
  4 & $1.845\times$ & 46.1\%  \\
  8 & $1.963\times$ & 24.5\%  \\
  \bottomrule
  \end{tabular}
  \caption{Pthread V5 平均加速比与并行效率（以 Pthread 1 线程为基准）}
  \label{tab:pthread-speedup}
\end{table}

Pthread V5 在所有线程数下都比 V2/V3/V4 更快，8 线程下平均约 $23.24$\,ms，
比单线程的 $45.6$\,ms 提升约 $2\times$。加速曲线与 OpenMP 趋势一致：
随着线程数增加，并行效率持续下降，主要受 bit-reversal 串行、stageRoots
串行预计算与 barrier 同步开销的影响。

%================================================================
\section{OpenMP 与 Pthread 对比}
%================================================================

\subsection{8 线程性能对比}

\begin{table}[H]
  \centering
  \begin{tabular}{lcccc}
  \toprule  
  \textbf{实现}   & $p=7340033$ & $p=104857601$ & $p=469762049$ & \textbf{三组平均}\\
  \midrule
  OpenMP 8 线程  & 21.3026 & 21.9200 & 22.1383 & \textbf{21.7870} \\
  Pthread 8 线程 & 22.5886 & 23.6726 & 23.4548 & 23.2387 \\
  \bottomrule
  \end{tabular}
  \caption{OpenMP 8 线程与 Pthread V5 8 线程三组输入对比（单位：ms）}
  \label{tab:omp-vs-pthread}
\end{table}

OpenMP 8 线程相对 Pthread V5 8 线程整体快约 $6.7\%$。

\subsection{对比分析}

\begin{itemize}
  \item \textbf{性能}：OpenMP 略优，主要得益于运行时较成熟的线程池、调度器
        和编译器辅助优化；
  \item \textbf{编程复杂度}：OpenMP 仅需在循环外加若干指导语句，
        Pthread 需要手写 \texttt{pthread\_create}、\texttt{pthread\_join}、
        \texttt{pthread\_barrier\_t}、\texttt{WorkerArg} 结构体以及阶段划分；
  \item \textbf{线程管理}：OpenMP 内置线程池，跨 \texttt{parallel} 区域复用线程；
        Pthread 需开发者自行实现 whole-pipeline 才能避免反复创建/销毁线程；
  \item \textbf{同步机制}：OpenMP 使用 \texttt{for}/\texttt{single} 的隐式 barrier，
        简单稳定；Pthread 需显式管理 barrier，粒度更灵活但也更容易引入死锁等问题；
  \item \textbf{调度灵活性}：OpenMP 提供 \texttt{schedule(static|dynamic|guided)}
        等简单切换；Pthread 想要等价效果需要自己实现任务队列；
  \item \textbf{可维护性}：OpenMP 几乎不改动原算法结构，Pthread 需要对算法
        做较深入的重构。
\end{itemize}

\textbf{结论}：在 NTT 这类``stage 间需 barrier、stage 内并行''的计算上，OpenMP 在性能、
开发成本和可维护性上整体优于手写 Pthread，与实验指导书的预期一致。Pthread 的价值
更多体现在需要精细化控制线程协作的场景，本实验中 whole-pipeline 线程池便是这种
精细控制的具体体现。

%================================================================
\section{大模数与 CRT 进阶优化}
%================================================================

\subsection{大模数为何需要特殊处理}

第四个测试输入使用模数 $p = 1337006139375617 \approx 2^{50}$。该模数远超 32 位上限，
单次模乘的中间结果可能达到 $2^{100}$，必须使用 \texttt{\_\_int128} 或更高精度类型。
\texttt{\_\_int128} 在 64 位平台上由两个寄存器组成，单次乘法和取模的开销显著高于 32 位。

\subsection{V7：大模数标量 NTT}

V7 是直接基于 \texttt{\_\_int128} 的标量 NTT：所有 butterfly 中的乘法和加减运算均在
64/128 位下完成。正确性通过 \texttt{fCheckLarge}。

\noindent\textbf{延迟：}\;$109.006$\,ms.

\subsection{V8：大模数 OpenMP NTT}

V8 直接在 V7 基础上套用 V1 的 OpenMP 并行模式，线程数为 8。
其延迟降至 $45.704$\,ms，加速比约 $2.39\times$。

\subsection{V9：5 模数 CRT 初版}

V9 引入 CRT：选用 5 个 32 位 NTT-friendly 模数分别完成 NTT，随后用 Garner 合并到
目标大模数。其优点是把昂贵的 \texttt{\_\_int128} 模乘换成多个 32 位模乘；
缺点是模数过多导致 NTT 工作量倍增，且单纯按模数级并行时，5 个模数无法与 8 个线程对齐。
最终延迟 $53.0959$\,ms，反而略慢于直接大模数 OpenMP。

\subsection{V10-1：4 模数 + Montgomery CRT}

V10-1 将 CRT 模数数量减少到 4 个，并在每个小模数 NTT 内部使用 Montgomery 表示，
降低单次模乘代价。延迟为 $48.2172$\,ms，相比 V9 有改善，但仍接近直接大模数 OpenMP，
未充分发挥 8 线程的潜力。

\subsection{V10-2：4 模数分组 CRT + Montgomery NTT + Fast Garner}

V10-2 是本文最终的进阶版本，针对 8 线程平台做了如下设计：

\begin{enumerate}
  \item \textbf{CRT\_MOD\_COUNT = 4}：使用 4 个互素 NTT-friendly 模数
        $\{469762049,\, 754974721,\, 880803841,\, 998244353\}$ 及其原根 $\{3, 11, 26, 3\}$，
        其乘积可以覆盖卷积单点最大值；
  \item \textbf{Grouped OpenMP 并行}：将 8 个 OpenMP 线程分成 4 组，每组 2 个线程
        协同处理一个 CRT 小模数；组内通过嵌套 \texttt{parallel}/外层任务分发实现，
        组间互不干扰；
  \item \textbf{Montgomery NTT}：每个 32 位小模数的 NTT 都使用 Montgomery reduction，
        模乘开销降到最低；
  \item \textbf{Fast Garner 合并}：预计算 6 个模逆元 $\mathrm{inv}_{01}, \mathrm{inv}_{02},
        \mathrm{inv}_{12}, \mathrm{inv}_{03}, \mathrm{inv}_{13}, \mathrm{inv}_{23}$ 以及
        合并系数 $\mathrm{coeff}_0$\,--\,$\mathrm{coeff}_3$，在合并阶段只剩若干次乘加。
\end{enumerate}

V10-2 在 8 线程下三次复跑结果：

\begin{table}[H]
  \centering
  \begin{tabular}{cc}
  \toprule  
  \textbf{Run} & \textbf{延迟 (ms)}\\
  \midrule
  Run 1 & 34.3048 \\
  Run 2 & 33.6297 \\
  Run 3 & \textbf{32.9417} \\
  \midrule
  平均  & \textbf{33.6254} \\
  \bottomrule
  \end{tabular}
  \caption{V10-2 三次独立运行的稳定性测试}
  \label{tab:v10-stability}
\end{table}

\subsection{大模数版本汇总对比}

\begin{table}[H]
  \centering
  \renewcommand{\arraystretch}{1.15}
  \begin{tabular}{lp{6.5cm}c}
  \toprule  
  \textbf{版本} & \textbf{说明} & \textbf{时间 (ms)}\\
  \midrule
  V7    & 直接大模数标量 NTT，\texttt{\_\_int128}              & 109.006 \\
  V8    & 直接大模数 OpenMP NTT，8 线程                        & 45.704  \\
  V9    & 5 模数 CRT 初版，8 线程                              & 53.0959 \\
  V10-1 & 4 模数 + Montgomery CRT，8 线程                      & 48.2172 \\
  V10-2 & 4 模数分组 CRT + Montgomery + Fast Garner，
          8 线程三次平均                                       & \textbf{33.6254} \\
  \bottomrule
  \end{tabular}
  \caption{大模数版本性能对比（$p = 1337006139375617$）}
  \label{tab:large-modulus}
\end{table}

\subsection{V10-2 为何最快}

V10-2 相对其它大模数实现的优势可归纳为以下四点：

\begin{enumerate}
  \item \textbf{昂贵 128 位乘法 $\to$ 廉价 32 位乘法}：相比 V7/V8，将单次大模数乘法
        替换为 4 个独立的 32 位 Montgomery 模乘，单次操作开销显著下降；
  \item \textbf{模数数量适配线程数}：相比 V9 的 5 模数方案，4 模数恰好可被 8 线程整除，
        每组 2 个线程；模数级并行与组内 stage 级并行结合，线程利用率高；
  \item \textbf{组内细粒度并行}：每个小模数 NTT 内部仍采用 grouped OpenMP 并行，
        避免了单纯按模数划分时``一个模数只用 1 个线程''带来的资源浪费；
  \item \textbf{Garner 合并开销降到最低}：通过预计算 $\mathrm{inv}_{ij}$ 和合并系数，
        将原本需要多次大数运算的合并阶段压缩到若干次乘加。
\end{enumerate}

\textbf{加速比小结}：

\begin{itemize}
  \item 相对 V7：$109.006 / 33.6254 \approx 3.24\times$
  \item 相对 V8：$45.704 / 33.6254 \approx 1.36\times$
  \item 相对 V9：$53.0959 / 33.6254 \approx 1.58\times$
\end{itemize}

%================================================================
\section{性能分析与讨论}
%================================================================

\subsection{正确性总结}

\begin{itemize}
  \item 前三个 32 位输入（$p = 7340033, 104857601, 469762049$）在所有版本（V0--V6）
        下均通过 \texttt{fCheck}；
  \item 第四个大模数输入（$p = 1337006139375617$）在 V7、V8、V9、V10-1、V10-2 下
        均通过 \texttt{fCheckLarge}；
  \item 在 $1/2/4/8$ 线程下重复实验，结果均稳定通过检查。
\end{itemize}

\subsection{扩展性分析}

OpenMP 与 Pthread V5 在 $1 \to 2 \to 4 \to 8$ 线程时均持续加速，但加速比从 2 线程的
$1.4\times$ 左右下降到 8 线程的 $2\times$ 左右，即并行效率从约 $70\%$ 降到约 $25\%$。
这一规律与 NTT 算法本身的特点高度一致：

\begin{enumerate}
  \item bit-reversal 阶段串行段占比固定，按 Amdahl 定律限制加速比；
  \item stageRoots 预计算由一个线程完成，每层都有一段串行段；
  \item 每个 stage 末尾的 barrier 在 8 线程下同步成本上升；
  \item 后期 stage 中 $b$ 增大、block 数减少，可并行任务量下降；
  \item NTT 数据规模为 $L = 2^{17}$ 时已超出 L1/L2 cache，访存成为瓶颈，
        线程数继续增加只会加重内存带宽压力。
\end{enumerate}

\subsection{Pthread 任务划分策略对比}

实验数据表明：

\begin{itemize}
  \item \textbf{V2 block-parallel} 在后期 stage 出现线程闲置；
  \item \textbf{V3 flat-butterfly} 解决了后期闲置，相对 V2 提升约 $10\%$；
  \item \textbf{V4 连续 flat} 因 NTT 访存模式不连续，没有进一步收益；
  \item \textbf{V5 whole-pipeline 线程池} 通过减少 \texttt{pthread\_create}/
        \texttt{pthread\_join} 次数，将延迟从 $35$\,ms 量级压到 $25$--$26$\,ms 量级，
        是 Pthread 路线中最显著的一次性能跃迁；
  \item \textbf{V6 cyclic-flat} 反而退化，表明``过度细化任务划分''并不总是好事，
        过于细的任务划分会引入整数除法和坏的访存局部性。
\end{itemize}

\subsection{大模数 CRT 收益分析}

CRT 路线对大模数 NTT 的收益主要来自算法层而非线程数堆叠。V8 与 V10-2 都使用 8 线程，
但 V10-2 通过``Montgomery 32 位模乘 + 4 模数分组 + Fast Garner''三件组合，
将平均延迟从 $45.7$\,ms 降到 $33.6$\,ms，再次说明在多线程已经接近平台极限时，
\textbf{算法级优化往往比单纯提高线程数更有效}。

\subsection{非线性加速的主要瓶颈}

综合 OpenMP、Pthread 各版本数据，本实验中无法获得线性加速的主要原因有：

\begin{itemize}
  \item bit-reversal 阶段串行部分；
  \item stageRoots 预计算阶段由单线程完成；
  \item 每层 NTT 后的 barrier 同步；
  \item 后期 stage 并行粒度变小；
  \item NTT 数据规模超出 cache，内存带宽成为瓶颈；
  \item CRT 路线中合并阶段虽然已优化，但仍有不可消除的常数开销。
\end{itemize}

%================================================================
\section{AI 辅助过程说明}
%================================================================

按照课程进阶要求允许并鼓励的方式，本实验在以下环节使用了生成式 AI 辅助：

\begin{itemize}
  \item \textbf{资料解读}：辅助梳理 Pthread/OpenMP 实验指导书与 NTT 选题介绍中的
        实验要求与评分细则；
  \item \textbf{并行化方案讨论}：辅助讨论 NTT stage 内 butterfly 是否独立、
        bit-reversal 是否需要串行、OpenMP \texttt{parallel}/\texttt{single}/\texttt{for}
        组合的语义边界；
  \item \textbf{版本演进与反思}：协助梳理从 V2 到 V6 的实验现象，包括``后期 stage 闲置''
        ``线程创建销毁开销''``cyclic 分配破坏局部性''等问题的分析路径；
  \item \textbf{CRT 思路细节}：辅助查阅与梳理 Garner 合并的数学推导，以及 4 模数
        grouped CRT 与 Montgomery NTT 的实现细节；
  \item \textbf{数据整理与报告组织}：协助整理实验数据、生成对比表格、组织报告章节顺序。
\end{itemize}

所有最终代码均经过本人理解、修改与本地编译运行；所有实验数据来源于本人在实验平台上的实测，
正确性通过框架提供的 \texttt{fCheck} 与 \texttt{fCheckLarge} 验证。
完整的 AI 辅助对话记录随报告附录或附件提交。

%================================================================
\section{总结与展望}
%================================================================

\subsection{主要工作与结论}

本文围绕 NTT 多项式乘法选题，完成了 Pthread 与 OpenMP 的多线程实现，并进一步实现了
基于 CRT 的大模数进阶优化。主要工作与结论可总结为：

\begin{enumerate}
  \item 完成 OpenMP NTT 实现（V1），在 32 位输入上 8 线程平均延迟 $21.7870$\,ms，
        相对 V0 标量基线加速约 $2.10\times$；
  \item 完成 5 个 Pthread 版本（V2--V6），其中 V5 whole-pipeline 线程池最优，
        8 线程平均 $23.2387$\,ms；
  \item 系统比较 OpenMP 与 Pthread 在性能、编程复杂度、同步机制、线程管理与
        调度灵活性上的差异，在 NTT 这类 stage 化的并行计算上，OpenMP 综合优于 Pthread；
  \item 完成大模数 NTT 的标量与 OpenMP 实现（V7、V8），分别将延迟从 $109.006$\,ms
        降到 $45.704$\,ms；
  \item 完成 CRT 进阶优化（V9、V10-1、V10-2），最终 V10-2 在 8 线程下三次平均
        $33.6254$\,ms，相对 V7 加速约 $3.24\times$，相对 V8 加速约 $1.36\times$。
\end{enumerate}

\subsection{经验总结}

\begin{itemize}
  \item 任务划分要与算法各 stage 的并行粒度匹配，过粗导致负载不均，过细引入额外
        计算与坏的局部性；
  \item 同步策略与线程管理开销在中小规模算法上往往是主要瓶颈，whole-pipeline 线程池
        在本实验中收益显著；
  \item 在多线程已经接近平台极限时，\textbf{算法级优化（如 Montgomery 模乘、CRT、
        Fast Garner）} 往往比单纯增加线程数更有效；
  \item OpenMP 的简洁性在工程中具有较高价值，Pthread 适合需要精细化控制的场景。
\end{itemize}

\subsection{展望}

\begin{itemize}
  \item 进一步探索 bit-reversal 阶段的并行化，减小 NTT 中的串行段；
  \item 实现全局共享的 stageRoots 预计算，使每层不再依赖``单线程预计算 + barrier''；
  \item 探索 NTT 选题介绍中提到的四分 NTT 算法，进一步降低 stage 数；
  \item 在不同的硬件平台（x86 / ARM Kunpeng）上重复实验，比较平台特性对策略选择的影响；
  \item 在更大规模（$L \ge 2^{20}$）的多项式乘法上检验各方案的扩展性。
\end{itemize}

%================================================================
\appendix
\section{关键代码片段}
%================================================================

\subsection{Montgomery 32 位模乘核心}

\begin{lstlisting}[title=Montgomery 模乘核心,frame=trbl,language={C++}]
struct Mont32 {
    uint32_t p;       // `模数`
    uint32_t pinv;    // -p^{-1} mod 2^32
    uint32_t r2;      // (2^32)^2 mod p
    uint32_t one;     // 2^32 mod p
};

// redc(T) = (T + (T*pinv mod 2^32) * p) >> 32
static inline uint32_t mont_redc(uint64_t T, const Mont32& m) {
    uint32_t k = (uint32_t)T * m.pinv;
    uint64_t t = T + (uint64_t)k * m.p;
    uint32_t r = (uint32_t)(t >> 32);
    return r >= m.p ? r - m.p : r;
}

static inline uint32_t mont_mul(uint32_t a, uint32_t b, const Mont32& m) {
    return mont_redc((uint64_t)a * b, m);
}
\end{lstlisting}

\subsection{V10-2 Grouped CRT 主体结构（节选）}

\begin{lstlisting}[title=4 模数 grouped CRT 主体,frame=trbl,language={C++}]
constexpr int CRT_MOD_COUNT = 4;
constexpr int THREAD_COUNT  = 8;
// `每组 2 个线程协同处理一个小模数`
constexpr int GROUP_THREADS = THREAD_COUNT / CRT_MOD_COUNT;

const uint32_t MODS[CRT_MOD_COUNT] = {
    469762049u, 754974721u, 880803841u, 998244353u
};
const uint32_t ROOTS[CRT_MOD_COUNT] = { 3u, 11u, 26u, 3u };

void crt_polynomial_multiply(/* ... */) {
    // `顶层：CRT_MOD_COUNT 个任务，每个任务用 GROUP_THREADS 个线程`
    #pragma omp parallel num_threads(CRT_MOD_COUNT)
    {
        int mod_id = omp_get_thread_num();
        // `在每个小模数内部，组内 2 个线程协同做 NTT`
        ntt_pipeline_for_modulus(mod_id, GROUP_THREADS);
    }

    // Garner `合并：使用预计算的 inv_ij 与 coeff_i`
    fast_garner_merge(/* ... */);
}
\end{lstlisting}

\subsection{Fast Garner 预计算}

\begin{lstlisting}[title=Fast Garner 预计算,frame=trbl,language={C++}]
struct FastGarner4 {
    uint32_t inv01, inv02, inv12, inv03, inv13, inv23;
    // `合并系数（在目标大模数下）`
    uint64_t coeff0, coeff1, coeff2, coeff3;
};

void prepare_fast_garner(FastGarner4& g) {
    g.inv01 = mod_inv(MODS[0], MODS[1]);
    g.inv02 = mod_inv(MODS[0], MODS[2]);
    g.inv12 = mod_inv(MODS[1], MODS[2]);
    g.inv03 = mod_inv(MODS[0], MODS[3]);
    g.inv13 = mod_inv(MODS[1], MODS[3]);
    g.inv23 = mod_inv(MODS[2], MODS[3]);
    // `合并到目标模数 P 的系数`
    g.coeff0 = 1ULL;
    g.coeff1 = (uint64_t)MODS[0] % P;
    g.coeff2 = mulmod_u64(g.coeff1, MODS[1], P);
    g.coeff3 = mulmod_u64(g.coeff2, MODS[2], P);
}
\end{lstlisting}

%----------------------------------------------------------------
\newpage
\begin{thebibliography}{9}

\bibitem{butenhof1997programming}
D. R. Butenhof.
\newblock \emph{Programming with POSIX Threads}.
\newblock Addison-Wesley, 1997.

\bibitem{dagum1998openmp}
L. Dagum and R. Menon.
\newblock OpenMP: An Industry-Standard API for Shared-Memory Programming.
\newblock \emph{IEEE Computational Science and Engineering}, 1998.

\bibitem{knuth1997art}
D. E. Knuth.
\newblock \emph{The Art of Computer Programming, Volume 2: Seminumerical Algorithms}.
\newblock Addison-Wesley, 1997.

\bibitem{cooley1965algorithm}
J. W. Cooley and J. W. Tukey.
\newblock An Algorithm for the Machine Calculation of Complex Fourier Series.
\newblock \emph{Mathematics of Computation}, 1965.

\bibitem{montgomery1985modular}
P. L. Montgomery.
\newblock Modular Multiplication Without Trial Division.
\newblock \emph{Mathematics of Computation}, 1985.

\bibitem{longa2016speeding}
P. Longa and M. Naehrig.
\newblock Speeding up the Number Theoretic Transform for Faster Ideal Lattice-Based Cryptography.
\newblock Cryptology ePrint Archive, 2016.

\bibitem{harvey2014faster}
D. Harvey.
\newblock Faster Arithmetic for Number-Theoretic Transforms.
\newblock Journal of Symbolic Computation, 2014.

\end{thebibliography}
\end{document}

